\section{Tarefa 5.2: Projeto e montagem da celula basica de um comparador}

Foi solicitado o projeto da celula basica de um comparador iterativo de duas palavras de $n$ bits. As saidas devem indicar, a cada etapa, se a palavra $A$ e maior que $B$, igual a $B$ ou menor que $B$. O relatorio deve apresentar a tabela-verdade da celula, as funcoes logicas obtidas, a montagem em protoboard e a descricao do funcionamento da rede iterativa.

\subsection{Tabela-verdade da celula}

\begin{table}[H]
    \centering
    \caption{Tabela-verdade da celula basica do comparador.}
    \begin{tabular}{ccccccc}
        \toprule
        $A_i$ & $B_i$ & $Y_{1,i}$ & $Y_{3,i}$ & $Y_{1,i-1}$ & $Y_{2,i-1}$ & $Y_{3,i-1}$ \\
        \midrule
        0 & 0 & 0 & 0 & & & \\
        0 & 0 & 0 & 1 & & & \\
        0 & 1 & 0 & 0 & & & \\
        0 & 1 & 0 & 1 & & & \\
        1 & 0 & 0 & 0 & & & \\
        1 & 0 & 0 & 1 & & & \\
        1 & 1 & 0 & 0 & & & \\
        1 & 1 & 0 & 1 & & & \\
        \bottomrule
    \end{tabular}
\end{table}

\subsection{Funcoes logicas}

\[
Y_{1,i-1} = \underline{\hspace{8cm}}
\]

\[
Y_{2,i-1} = \underline{\hspace{8cm}}
\]

\[
Y_{3,i-1} = \underline{\hspace{8cm}}
\]

\subsection{Montagem}

\begin{figure}[H]
    \centering
    % \includegraphics[width=0.8\textwidth]{questao2/comparador_1bit.png}
    \fbox{\rule{0pt}{5cm}\rule{0.8\textwidth}{0pt}}
    \caption{Espaco reservado para a montagem da celula do comparador.}
\end{figure}

\subsection{Resultados experimentais}

\begin{table}[H]
    \centering
    \caption{Resultados obtidos para a celula basica do comparador.}
    \begin{tabular}{ccccccc}
        \toprule
        $A_i$ & $B_i$ & $Y_{1,i}$ & $Y_{3,i}$ & $Y_{1,i-1}$ & $Y_{2,i-1}$ & $Y_{3,i-1}$ \\
        \midrule
        & & & & & & \\
        & & & & & & \\
        & & & & & & \\
        & & & & & & \\
        \bottomrule
    \end{tabular}
\end{table}

\subsection{Rede iterativa}

Explique aqui como as celulas sao encadeadas para formar um comparador de $n$ bits e como se obtem as saidas globais $Y_1$, $Y_2$ e $Y_3$.
